\documentclass[twoside]{article}
\usepackage[utf8]{inputenc}

% The package supports other formats as well, but this template
% is primarily checked and updated for examinations at TUM BGDM
%\usepackage[uebung, answers]{tumbgdm}

\usepackage[uebung]{tumbgdm}
\usepackage{pythonhighlight}

\Professor{Prof. Dr. Martin Werner}
\Universitaet{Technical University of Munich}
\Institut{School of Engineering and Design}
\Lehrstuhl{Professorship Big Geospatial Data Management}
\Semester[WiTe 2024/25]{Winter Term 2024/25}
\Vorlesung{Principles of Programming}
\InstitutsWebsite{https://www.bgd.lrg.tum.de/}
\VorlesungsWebsite{https://www.moodle.tum.de/course/view.php?id=103704}

\nummerUebungsblatt{2}
\nummerErsteFrage{1}

\LearningOutcome{ Loops, Conditions, Lists, List Comprehensions}

\begin{document}
\maketitle
\vspace{10pt}
\\

\section{Python Loops}

\begin{task}{For Loops}{}{}
Use a for loop to print a multiplication table from 1 to 10.

\vspace{10pt}
Expected output:
\begin{lstlisting}[numbers=none]
1   2  3  4  5  6  7  8  9  10 		
2   4  6  8 10 12 14 16 18  20 		
3   6  9 12 15 18 21 24 27  30 		
4   8 12 16 20 24 28 32 36  40 		
5  10 15 20 25 30 35 40 45  50 		
6  12 18 24 30 36 42 48 54  60 		
7  14 21 28 35 42 49 56 63  70 		
8  16 24 32 40 48 56 64 72  80 		
9  18 27 36 45 54 63 72 81  90 		
10 20 30 40 50 60 70 80 90 100 
\end{lstlisting}

\end{task}

\begin{task}{While Loops}{}{}
Using a while statement, write a function that takes two int values ("base" and "exponent") as input and prints the base to the power of the exponent.

\vspace{10pt}
\textbf{Example}\\
Expected output for input base=2 and exp=5:
\begin{lstlisting}[numbers=none]
2 to the power of 5 is:  32
\end{lstlisting}

\end{task}

\clearpage
\section{Python List Comprehensions}

Solve the following tasks with list comprehensions:
\begin{lstlisting}[numbers=none]
newlist = [expression for item in iterable if condition] 
\end{lstlisting}


\begin{task}{Reverse String}{}{}
Write a function that reverses a string.

\vspace{10pt}
\textbf{Example}\\
Expected output for input s = "hello":
\begin{lstlisting}[numbers=none]
["o","l","l","e","h"]
\end{lstlisting}

%\vspace{5pt}
%\textbf{Constraints:}
%\begin{itemize}
%    \item $1 \leq len(s) \leq 10^5$
%    \item s[i] is a printable ascii character.
%\end{itemize}
\end{task}


\begin{task}{Sort Array}{}{}
Given an integer array nums, move all the even integers at the beginning of the array followed by all the odd integers.
Return any array that satisfies this condition.

\vspace{10pt}
\textbf{Example}\\
Expected output for input nums = [3,1,2,4]:
\begin{lstlisting}[numbers=none]
[2,4,3,1]
\end{lstlisting}
Explanation: The outputs [4,2,3,1], [2,4,1,3], and [4,2,1,3] would also be accepted.

%\textbf{Constraints:}
%\begin{itemize}
%    \item $1 \leq len(nums) \leq 5000$
%    \item $0 \leq nums[i] \leq 5000$
%\end{itemize}
\end{task}


\begin{task}{Fizz Buzz}{}{}
Fizz buzz is a group word game for children to teach them about division. Players take turns to count incrementally, replacing any number divisible by three with the word "fizz", and any number divisible by five with the word "buzz", and any number divisible by both 3 and 5 with the word "fizzbuzz" (from \href{https://en.wikipedia.org/wiki/Fizz_buzz}{Wikipedia}).\\

Given an integer n, return a string array answer (1-indexed) where:

\begin{enumerate}
    \item answer[i] == “FizzBuzz” if i is divisible by 3 and 5.
    \item answer[i] == “Fizz” if i is divisible by 3.
    \item answer[i] == “Buzz” if i is divisible by 5.
    \item answer[i] == i (as a string) if none of the above conditions are true.
\end{enumerate}

\vspace{10pt}
\textbf{Example}\\
Expected output for input n = 15:
\begin{lstlisting}[numbers=none]
["1","2","Fizz","4","Buzz","Fizz","7","8","Fizz","Buzz","11","Fizz","13","14","FizzBuzz"]
\end{lstlisting}

%\vspace{5pt}
%\textbf{Constraints:}
%\begin{itemize}
%    \item $1 \leq n \leq 10^4$
%\end{itemize}
\end{task}


\end{document}